% Part: second-order-logic
% Chapter: syntax-and-semantics
% Section: introduction

\documentclass[../../../include/open-logic-section]{subfiles}

\begin{document}

% SEGMENT OLP-0324-S01

\olfileid{sol}{syn}{int}

\olsection{அறிமுகம்}

முதல்தரத் தருக்கத்தில், கொடுக்கப்பட்ட ஒரு மொழியின் தருக்கமல்லாத
குறிகளை, அதாவது அதன் !!{constant}s, !!{function}s, !!{predicate}s
ஆகியவற்றைத் தருக்கக் குறிகளுடன் சேர்த்து, முதல்தர
!!{structure}s பற்றியவற்றை வெளிப்படுத்துகிறோம். நிறைவுறுத்தல் என்ற
கருத்தைப் பயன்படுத்தி இதைச் செய்கிறோம். அது !!a{structure}~$\Struct{M}$,
ஒரு மாறி மதிப்பீடு~$s$, !!a{formula}~$!A$ ஆகியவற்றைத் தொடர்புபடுத்துகிறது:
$\Sat{M}{!A}[s]$ பொருந்துவது அப்போதும் அப்போது மட்டுமே, $!A$ இன்
!!{constant}s, !!{function}s, !!{predicate}s ஆகியவற்றை
$\Struct{M}$ கூறுவதுபோலவும், அதன் கட்டற்ற மாறிகளை~$s$ கூறுவதுபோலவும்
பொருள்கொண்டால் அது வெளிப்படுத்துவது மெய்யாகும். !!{identity}~$\eq$ இன்
பொருள்கோள் $\Sat{M}{!A}[s]$ இன் வரையறையிலேயே உள்ளமைக்கப்பட்டுள்ளது;
அவ்வாறே~$\lforall$, $\lexists$ ஆகியவற்றின் பொருள்கோளும்
உள்ளமைக்கப்பட்டுள்ளது. முதலாவது குறி எப்போதும் அந்த அமைப்பின்
!!{domain}~$\Domain{M}$ மீதான சமனுறவாகப் பொருள்கொள்ளப்படுகிறது;
அளவையடைகள் எப்போதும் முழு !!{domain} மீதும் வீச்சுக்கொண்டவையாகப்
பொருள்கொள்ளப்படுகின்றன. ஆனால், முக்கியமாக, !!{domain} இன் உறுப்புகள்
மீது மட்டுமே அளவையிடல் அனுமதிக்கப்படுகிறது; எனவே ஓர் அளவையடைக்குப்
பின்னர் பொருள் !!{variable}s மட்டுமே வரலாம்.

இரண்டாம்தரத் தருக்கத்தில், கட்டற்ற மற்றும் கட்டுண்ட சார்பு மாறிகளையும்
பயனிலை மாறிகளையும், அவற்றின்மீதான அளவையிடலையும் சேர்க்கும் வகையில்
மொழியும் நிறைவுறுத்தலின் வரையறையும் நீட்டிக்கப்படுகின்றன. பொருள்
மாறிகள் !!{constant}s உடன் தொடர்புடைய அதே முறையில் இந்த மாறிகள்
!!{function}s மற்றும் !!{predicate}s உடன் தொடர்புடையவை. இரண்டாம்தரத்
தருக்கத்தின் உறுப்புகளையும் !!{formula}s ஐயும் அமைப்பதில் அவை அதே
பங்கை வகிக்கின்றன; அவற்றின்மீதான அளவையிடலும் ஒத்த முறையில்
கையாளப்படுகிறது. \emph{தரநிலை} பொருள்கோளில், இரண்டாம்தர அளவையடைகள்
சரியான வகையைச் சேர்ந்த அனைத்து சாத்தியமான பொருள்கள்மீதும்
வீச்சுக்கொள்கின்றன (சார்பு மாறிகளுக்கு $\Domain{M}$ இலிருந்து
~$\Domain{M}$ க்குச் செல்லும் $n$-இடச் சார்புகளும், பயனிலை
மாறிகளுக்கு $n$-இட உறவுகளும்). எடுத்துக்காட்டாக,
$\lforall[\Obj{v_0}][(\Obj{P^1_0}(\Obj{v_0}) \lor \lnot
  \Obj{P^1_0}(\Obj{v_0}))]$ என்பது முதல்தர மற்றும் இரண்டாம்தரத்
தருக்கங்கள் இரண்டிலும் ஒரு வாய்பாடு. இரண்டாம்தரத் தருக்கத்தில்
$\lforall[\Obj{V^1_0}][\lforall[\Obj{v_0}][(\Obj{V^1_0}(\Obj{v_0}) \lor
    \lnot \Obj{V^1_0}(\Obj{v_0}))]]$ மற்றும்
$\lexists[\Obj{V^1_0}][\lforall[\Obj{v_0}][(\Obj{V^1_0}(\Obj{v_0}) \lor
    \lnot \Obj{V^1_0}(\Obj{v_0}))]]$ ஆகியவற்றையும் கருதலாம்.
இவற்றில் கட்டற்ற மாறிகள் இல்லாததால், இவை இரண்டாம்தரத் தருக்கத்தின்
!!{sentence}s. இங்கு $\Obj{V^1_0}$ என்பது $1$-இட இரண்டாம்தரப் பயனிலை
மாறி. $\Obj{V^1_0}$ க்கு அனுமதிக்கப்படும் பொருள்கோள்கள்,
$\Obj{P^1_0}$ போன்ற $1$-இட !!{predicate} ஒன்றுக்கு நாம் அளிக்கக்கூடிய
பொருள்கோள்களே, அதாவது~$\Domain{M}$ இன் உட்கணங்கள். அவற்றின்மீதான
அளவையிடல் என்பது
$\lforall[\Obj{v_0}][(\Obj{V^1_0}(\Obj v_0) \lor \lnot
  \Obj{V^1_0}(v_0))]$ ஆனது $\Obj{V^1_0}$ இன் மதிப்பாக
~$\Domain{M}$ இன் ஓர் உட்கணத்தை அளிக்கும் எல்லா வழிகளிலும், அல்லது
குறைந்தது ஒரு வழியிலாவது, பொருந்துகிறது என்று கூறுவதாகும். எந்தக்
கணமும் ஒரு குறிப்பிட்ட பொருளைக் கொண்டிருக்கும் அல்லது கொண்டிராது
என்பதால், இரு வாக்கியங்களும் எந்த !!{structure} இலும் மெய்யாகும்.
\end{document}
